%% Dans ce fichier, les �ventuels accents sont cod�s
%% en ISO-8859-15 (aka latin9)
%% La ligne suivante donne trois exemples 
%% e accent aigu : �, e accent grave : �, e-dans-l'o : �
%%%%%ENTREES%%%%
\begin{tabarg}
\multicolumn{3}{l}{\titretab {Entr\'ees}}\\
\hline
$X$ & Vecteur, \texttt{float[$N_{level}$]} & Signal lidar \\
$S_1$ & Coeff., \texttt{int} & Seuil de d�tection \\
$S_2$ & Coeff., \texttt{int} & Seuil de d�tection \\
$Z$ & Vecteur, \texttt{float[$N_{level}$]} & Altitudes r�elles des portes du lidar (m)\\
$Z_{min}$ & Coeff., \texttt{int} & Altitude en dessous de laquelle on ne d�tecte pas de nuage (m)\\
\end{tabarg}

\medskip
%%%%%SORTIES%%%%
\begin{tabarg}
\multicolumn{3}{l}{\titretab {Sorties}}\\
\hline
$Base$ & Vecteur, \texttt{float[6]} & Altitude des bases des nuages (m)\\
$Max$ & Vecteur, \texttt{float[6]} & Altitude des maximums de r�tro-diffusion (m)\\
$Sommet$ & Vecteur, \texttt{float[6]} & Altitude des sommets des nuages (m)\\
$Flag_{nuage}$ & Vecteur, \texttt{float} & Nombre de couches d�tect�es (m)\\
\end{tabarg}

